\documentclass[12pt, a4paper]{article}

% ── Paquetes ──────────────────────────────────────────────────────────────────
\usepackage[utf8]{inputenc}
\usepackage[T1]{fontenc}
\usepackage[spanish]{babel}
\usepackage{amsmath, amssymb}
\usepackage{booktabs}
\usepackage{array}
\usepackage{geometry}
\usepackage{hyperref}
\usepackage{listings}
\usepackage{xcolor}
\usepackage{float}
\usepackage{lmodern}

\geometry{margin=2.5cm}

\hypersetup{colorlinks=true, linkcolor=blue, urlcolor=blue, citecolor=blue}

% ── Estilo para código Python ──────────────────────────────────────────────────
\lstdefinestyle{python}{
  language=Python,
  basicstyle=\ttfamily\footnotesize,
  keywordstyle=\color{blue}\bfseries,
  stringstyle=\color{orange},
  commentstyle=\color{gray}\itshape,
  numberstyle=\tiny\color{gray},
  numbers=left,
  stepnumber=1,
  breaklines=true,
  frame=single,
  backgroundcolor=\color{gray!8},
  showstringspaces=false,
  tabsize=4
}

\lstset{style=python}

% ── Portada ───────────────────────────────────────────────────────────────────
\title{%
  \LARGE \textbf{Series de Tiempo -- Notebook de Prácticas} \\[0.5em]
  \large Finanzas Cuantitativas
}
\author{}
\date{}

% ══════════════════════════════════════════════════════════════════════════════
\begin{document}

\maketitle
\tableofcontents
\newpage

% ══════════════════════════════════════════════════════════════════════════════
\section{Carga de Datos}

\begin{lstlisting}
import pandas as pd

ventas = pd.read_csv("california_housing_train.csv")
ventas.head(51)
\end{lstlisting}

% ══════════════════════════════════════════════════════════════════════════════
\section{Ajustar y Pronosticar con SARIMA (Modelo Airline)}

\subsection{Ajuste del modelo SARIMA$(0,1,1)(0,1,1)_{12}$ a precios de AMZN}

Se importa la clase \texttt{SARIMAX} de \texttt{statsmodels} y se define el orden
no estacional $(0,1,1)$ y estacional $(0,1,1,12)$:

\begin{lstlisting}
import pandas as pd
import numpy as np
from statsmodels.tsa.statespace.sarimax import SARIMAX
import yfinance as yf

# Descargar datos de AMZN
amzn = yf.download("AMZN", end="2026-01-01", progress=False)
precios_amzn = amzn["Close"]

# Establecer frecuencia de dias habiles
precios_amzn = precios_amzn.asfreq('B')

# Definir el orden del modelo SARIMA
order = (0, 1, 1)
seasonal_order = (0, 1, 1, 12)

print(f"Ajustando SARIMA{order}{seasonal_order}...")

sarima_model = SARIMAX(precios_amzn, order=order,
                       seasonal_order=seasonal_order,
                       enforce_stationarity=False,
                       enforce_invertibility=False)

sarima_results = sarima_model.fit()
print(sarima_results.summary())
\end{lstlisting}

\subsection{Generación de pronósticos con intervalos de confianza}

\begin{lstlisting}
import matplotlib.pyplot as plt

forecast_steps_sarima = 30

forecast_sarima_output = sarima_results.get_forecast(
    steps=forecast_steps_sarima)
forecast_sarima_mean = forecast_sarima_output.predicted_mean
conf_int_sarima = forecast_sarima_output.conf_int(alpha=0.05)

plt.figure(figsize=(15, 7))
plt.plot(precios_amzn.index, precios_amzn,
         label='Serie Original (AMZN Close)', color='blue')
plt.plot(forecast_sarima_mean.index, forecast_sarima_mean,
         label='Pronostico SARIMA', color='red', linestyle='--')
plt.fill_between(conf_int_sarima.index,
                 conf_int_sarima.iloc[:, 0],
                 conf_int_sarima.iloc[:, 1],
                 color='red', alpha=0.2,
                 label='95% Intervalo de Confianza')
plt.title('Pronostico SARIMA(0,1,1)(0,1,1)12 de Precios AMZN')
plt.xlabel('Fecha')
plt.ylabel('Precio de Cierre')
plt.legend()
plt.grid(True)
plt.show()
\end{lstlisting}

\subsection{Resumen: Metodología SARIMA}

\paragraph{Orden no estacional $(p,d,q)=(0,1,1)$:}
\begin{itemize}
  \item $p=0$: Sin componente autorregresivo no estacional.
  \item $d=1$: Se aplica una diferencia ordinaria para lograr estacionariedad.
  \item $q=1$: Componente de media móvil de primer orden.
\end{itemize}

\paragraph{Orden estacional $(P,D,Q,s)=(0,1,1,12)$:}
\begin{itemize}
  \item $P=0$: Sin componente autorregresivo estacional.
  \item $D=1$: Diferenciación estacional anual.
  \item $Q=1$: Media móvil estacional de primer orden.
  \item $s=12$: Periodicidad mensual (ciclo anual).
\end{itemize}

Los intervalos de confianza del 95\% se amplían conforme aumenta el horizonte de
pronóstico, reflejando la creciente incertidumbre en predicciones de largo plazo.

% ══════════════════════════════════════════════════════════════════════════════
\section{Simulación de Procesos Estocásticos}

\subsection{Ruido blanco}

\begin{lstlisting}
import numpy as np
import matplotlib.pyplot as plt

np.random.seed(42)
white_noise = np.random.normal(loc=0, scale=1, size=200)

time = np.arange(len(white_noise))

plt.figure(figsize=(12, 6))
plt.plot(time, white_noise)
plt.title('Simulated Stationary White Noise Process')
plt.xlabel('Time')
plt.ylabel('Value')
plt.grid(True)
plt.show()
\end{lstlisting}

El ruido blanco es el ejemplo canónico de proceso estacionario: media y varianza
constantes en el tiempo, y autocovarianza nula para todo rezago $k \neq 0$.

\subsection{Proceso AR(1)}

La fórmula matemática del proceso AR(1) es:
\[
  X_t = c + \phi X_{t-1} + \varepsilon_t
\]
donde $|\phi| < 1$ garantiza la estacionariedad.

\begin{lstlisting}
import numpy as np
import matplotlib.pyplot as plt

np.random.seed(42)
white_noise = np.random.normal(loc=0, scale=1, size=200)

phi = 0.7
c = 0
ar1_process = np.zeros_like(white_noise)
ar1_process[0] = 0

for t in range(1, len(white_noise)):
    ar1_process[t] = c + phi * ar1_process[t-1] + white_noise[t]

plt.figure(figsize=(12, 6))
plt.plot(np.arange(len(ar1_process)), ar1_process)
plt.title(r'Proceso AR(1) Simulado ($\phi=0.7$)')
plt.xlabel('Tiempo')
plt.ylabel('Valor')
plt.grid(True)
plt.show()
\end{lstlisting}

% ══════════════════════════════════════════════════════════════════════════════
\section{Selección de Modelos ARIMA con AIC y BIC}

\subsection{Ajuste de múltiples modelos ARIMA}

\begin{lstlisting}
import pandas as pd
from statsmodels.tsa.arima.model import ARIMA

model_orders = {
    "AR(0)":     (0, 0, 0),
    "AR(1)":     (1, 0, 0),
    "MA(1)":     (0, 0, 1),
    "ARMA(1,1)": (1, 0, 1)
}

model_results = []

for model_name, order in model_orders.items():
    try:
        model   = ARIMA(ar1_process, order=order)
        results = model.fit()
        model_results.append({
            "Model": model_name, "Order": order,
            "AIC": results.aic, "BIC": results.bic
        })
        print(f"{model_name}: AIC={results.aic:.2f}, BIC={results.bic:.2f}")
    except Exception as e:
        print(f"No se pudo ajustar {model_name}: {e}")

model_results_df = pd.DataFrame(model_results)
display(model_results_df)
\end{lstlisting}

\subsection{Resultados y comparación}

\begin{table}[H]
  \centering
  \caption{Criterios de información para modelos ARIMA ajustados al proceso AR(1)}
  \label{tab:aic_bic}
  \begin{tabular}{llrr}
    \toprule
    \textbf{Modelo} & \textbf{Orden} & \textbf{AIC} & \textbf{BIC} \\
    \midrule
    AR(0)     & (0, 0, 0) & 648.66 & 655.26 \\
    \textbf{AR(1)} & \textbf{(1, 0, 0)} & \textbf{543.25} & \textbf{553.15} \\
    MA(1)     & (0, 0, 1) & 571.39 & 581.28 \\
    ARMA(1,1) & (1, 0, 1) & 545.25 & 558.44 \\
    \bottomrule
  \end{tabular}
\end{table}

El modelo \textbf{AR(1)} minimiza tanto el AIC como el BIC, identificando
correctamente la estructura verdadera del proceso simulado.

\subsection{Visualización de AIC y BIC}

\begin{lstlisting}
import matplotlib.pyplot as plt
import seaborn as sns

model_results_melted = model_results_df.melt(
    id_vars=['Model', 'Order'],
    value_vars=['AIC', 'BIC'],
    var_name='Metric', value_name='Value')

plt.figure(figsize=(12, 7))
sns.barplot(data=model_results_melted, x='Model', y='Value',
            hue='Metric', palette={'AIC': 'skyblue', 'BIC': 'lightcoral'})
plt.title('Comparacion de Modelos ARIMA por AIC y BIC')
plt.xlabel('Modelo')
plt.ylabel('Valor')
plt.legend(title='Metrica')
plt.grid(axis='y', linestyle='--', alpha=0.7)
plt.tight_layout()
plt.show()
\end{lstlisting}

% ══════════════════════════════════════════════════════════════════════════════
\section{Proceso MA(1): Simulación e Identificación}

\subsection{Fórmula matemática}

El proceso MA(1) se define como:
\[
  Y_t = \mu + \varepsilon_t + \theta_1\,\varepsilon_{t-1}
\]
donde $\mu$ es la media constante y $\theta_1$ es el coeficiente de media móvil.

\subsection{Simulación}

\begin{lstlisting}
import numpy as np

np.random.seed(42)
white_noise_ma = np.random.normal(loc=0, scale=1, size=200)

mean    = 0
theta_1 = 0.7
ma1_process = np.zeros_like(white_noise_ma)
ma1_process[0] = mean + white_noise_ma[0]

for t in range(1, len(white_noise_ma)):
    ma1_process[t] = mean + white_noise_ma[t] + theta_1 * white_noise_ma[t-1]
\end{lstlisting}

\subsection{Gráficas ACF y PACF}

\begin{lstlisting}
from statsmodels.graphics.tsaplots import plot_acf, plot_pacf
import matplotlib.pyplot as plt

fig, axes = plt.subplots(1, 2, figsize=(15, 6))

plot_acf(ma1_process, lags=20, ax=axes[0],
         title='ACF del Proceso MA(1)')
axes[0].set_xlabel('Rezago')
axes[0].set_ylabel('Autocorrelacion')
axes[0].grid(True)

plot_pacf(ma1_process, lags=20, ax=axes[1],
          title='PACF del Proceso MA(1)')
axes[1].set_xlabel('Rezago')
axes[1].set_ylabel('Autocorrelacion Parcial')
axes[1].grid(True)

plt.tight_layout()
plt.show()
\end{lstlisting}

\paragraph{Interpretación:}
\begin{itemize}
  \item \textbf{ACF}: presenta un pico significativo en el rezago 1 y luego
        se trunca bruscamente, identificando $q=1$.
  \item \textbf{PACF}: muestra un decaimiento gradual, consistente con un
        proceso MA.
\end{itemize}

% ══════════════════════════════════════════════════════════════════════════════
\section{Anotaciones en ACF y PACF para Identificación de Órdenes}

\begin{lstlisting}
from statsmodels.graphics.tsaplots import plot_acf, plot_pacf
import matplotlib.pyplot as plt

# ACF del MA(1): destacar q=1
fig, ax = plt.subplots(figsize=(10, 6))
plot_acf(ma1_process, lags=10, ax=ax,
         title='ACF del Proceso MA(1) (Destacando orden q)')
ax.axvline(x=1, color='red', linestyle='--', linewidth=2,
           label='Rezago significativo (q=1)')
ax.set_xlabel('Rezago')
ax.set_ylabel('Autocorrelacion')
ax.legend()
ax.grid(True)
plt.tight_layout()
plt.show()

# PACF del AR(1): destacar p=1
fig, ax = plt.subplots(figsize=(10, 6))
plot_pacf(ar1_process, lags=10, ax=ax,
          title='PACF del Proceso AR(1) (Destacando orden p)')
ax.axvline(x=1, color='red', linestyle='--', linewidth=2,
           label='Rezago significativo (p=1)')
ax.set_xlabel('Rezago')
ax.set_ylabel('Autocorrelacion Parcial')
ax.legend()
ax.grid(True)
plt.tight_layout()
plt.show()
\end{lstlisting}

\paragraph{Reglas de identificación:}
\begin{itemize}
  \item Para identificar $q$ (orden MA): buscar el corte abrupto en la
        \textbf{ACF}.
  \item Para identificar $p$ (orden AR): buscar el corte abrupto en la
        \textbf{PACF}.
  \item Si ambas funciones decaen gradualmente, se sugiere un proceso
        \textbf{ARMA} mixto.
\end{itemize}

% ══════════════════════════════════════════════════════════════════════════════
\section{Pronóstico del Proceso AR(1) con Intervalos de Confianza}

\begin{lstlisting}
from statsmodels.tsa.arima.model import ARIMA
import matplotlib.pyplot as plt
import numpy as np

# Ajustar ARIMA(1,0,0)
ar_model         = ARIMA(ar1_process, order=(1, 0, 0))
ar_model_results = ar_model.fit()

forecast_steps  = 20
forecast_output = ar_model_results.get_forecast(steps=forecast_steps)
forecast_mean   = forecast_output.predicted_mean
conf_int        = forecast_output.conf_int(alpha=0.05)

time_original = np.arange(len(ar1_process))
time_forecast = np.arange(len(ar1_process),
                           len(ar1_process) + forecast_steps)

plt.figure(figsize=(14, 7))
plt.plot(time_original, ar1_process,
         label='Proceso AR(1) Original', color='blue')
plt.plot(time_forecast, forecast_mean,
         label='Valores Pronosticados', color='red', linestyle='--')
plt.fill_between(time_forecast, conf_int[:, 0], conf_int[:, 1],
                 color='red', alpha=0.2,
                 label='95% Intervalo de Confianza')
plt.title('Pronostico ARIMA del Proceso AR(1)')
plt.xlabel('Tiempo')
plt.ylabel('Valor')
plt.legend()
plt.grid(True)
plt.show()
\end{lstlisting}

Los valores pronosticados muestran convergencia a la media de largo plazo
($c=0$), comportamiento esperado de un proceso AR(1) estacionario. Los
intervalos de confianza se ensanchan progresivamente con el horizonte.

% ══════════════════════════════════════════════════════════════════════════════
\section{Análisis de Datos Reales: Precios de Cierre de AMZN}

\subsection{Preparación de datos}

\begin{lstlisting}
import pandas as pd
import yfinance as yf

amzn = yf.download("AMZN", end="2026-01-01")
precios_amzn = amzn["Close"]

# Verificar tipo de indice
is_datetime = pd.api.types.is_datetime64_any_dtype(precios_amzn.index)
print(f"Indice datetime: {is_datetime}")

# Verificar y tratar valores faltantes
missing_total = precios_amzn.isnull().sum().item()
if missing_total > 0:
    precios_amzn = precios_amzn.ffill()

# Asignar frecuencia
precios_amzn = precios_amzn.asfreq('B')
print(precios_amzn.head(5))
\end{lstlisting}

\subsection{Selección de modelos ARIMA en datos reales}

\begin{lstlisting}
from statsmodels.tsa.arima.model import ARIMA
import pandas as pd

model_orders_real = {
    "ARIMA(0,1,0)": (0, 1, 0),
    "ARIMA(1,1,0)": (1, 1, 0),
    "ARIMA(0,1,1)": (0, 1, 1),
    "ARIMA(1,1,1)": (1, 1, 1)
}

model_results_real = []

for model_name, order in model_orders_real.items():
    try:
        model   = ARIMA(precios_amzn, order=order,
                        enforce_stationarity=False,
                        enforce_invertibility=False)
        results = model.fit()
        model_results_real.append({
            "Model": model_name, "Order": order,
            "AIC": results.aic, "BIC": results.bic
        })
        print(f"{model_name}: AIC={results.aic:.2f}, BIC={results.bic:.2f}")
    except Exception as e:
        print(f"Error en {model_name}: {e}")

model_results_real_df = pd.DataFrame(model_results_real)
display(model_results_real_df)
\end{lstlisting}

\begin{table}[H]
  \centering
  \caption{Criterios de información para modelos ARIMA ajustados a precios AMZN}
  \label{tab:aic_bic_real}
  \begin{tabular}{llrr}
    \toprule
    \textbf{Modelo} & \textbf{Orden} & \textbf{AIC} & \textbf{BIC} \\
    \midrule
    ARIMA(0,1,0)          & (0, 1, 0) & 95.41 & 96.46 \\
    ARIMA(1,1,0)          & (1, 1, 0) & 97.00 & 99.09 \\
    \textbf{ARIMA(0,1,1)} & \textbf{(0, 1, 1)} & \textbf{92.61} & \textbf{94.61} \\
    ARIMA(1,1,1)          & (1, 1, 1) & 94.61 & 97.60 \\
    \bottomrule
  \end{tabular}
\end{table}

El modelo \textbf{ARIMA$(0,1,1)$} presenta los menores valores de AIC y BIC.

\subsection{Pronóstico con Meta Prophet}

\begin{lstlisting}
import sys
# Instalar si es necesario:
# !{sys.executable} -m pip install prophet

import pandas as pd
from prophet import Prophet

# Preparar datos para Prophet
prophet_df = precios_amzn.reset_index()
prophet_df.columns = ['ds', 'y']
prophet_df['ds'] = pd.to_datetime(prophet_df['ds'])

# Ajustar modelo Prophet con intervalo de 95%
model = Prophet(interval_width=0.95)
model.fit(prophet_df)

# Crear dataframe futuro (30 dias habiles)
future   = model.make_future_dataframe(periods=30, freq='B')
forecast = model.predict(future)

print(forecast[['ds', 'yhat', 'yhat_lower', 'yhat_upper']].tail(10))
\end{lstlisting}

\begin{lstlisting}
import matplotlib.pyplot as plt

# Grafico de pronostico
fig1 = model.plot(forecast)
plt.title('Pronostico Prophet de Precios AMZN')
plt.xlabel('Fecha')
plt.ylabel('Precio de Cierre')
plt.grid(True)
plt.show()

# Componentes del pronostico (tendencia, estacionalidad)
fig2 = model.plot_components(forecast)
plt.show()
\end{lstlisting}

\paragraph{Interpretación:}
\begin{itemize}
  \item \texttt{yhat}: predicción puntual del precio de cierre.
  \item \texttt{yhat\_lower}/\texttt{yhat\_upper}: límites inferior y superior
        del intervalo de confianza al 95\%.
  \item Los intervalos se amplían con el horizonte de pronóstico, reflejando
        mayor incertidumbre en el largo plazo.
\end{itemize}

% ══════════════════════════════════════════════════════════════════════════════
\section{Preguntas de Aplicación}

\subsection{Pregunta 5 -- Serie S\&P 500 (GSPC)}

\begin{lstlisting}
import yfinance as yf
import pandas as pd

ticker = "^GSPC"
data   = yf.download(ticker, start="2014-01-01", end="2026-01-01")

# Serie mensual (media del mes)
ts = data['Close'].resample('MS').mean()
ts = ts.dropna()
\end{lstlisting}

\paragraph{b) Transformaciones sugeridas:}
\begin{itemize}
  \item \textbf{Frecuencia}: mensual (\texttt{MS}).
  \item \textbf{Valores faltantes}: interpolación lineal o eliminación de
        periodos vacíos.
  \item \textbf{Logaritmos}: aplicar $\ln(X_t)$ si la serie exhibe crecimiento
        exponencial.
  \item \textbf{Diferenciación}: identificar $d$ y $D$ mediante la prueba ADF.
\end{itemize}

\paragraph{c) Valores plausibles para $d$, $D$ y $s$:}
\begin{itemize}
  \item $s = 12$: ciclo estacional anual (datos mensuales).
  \item $d = 1$: necesaria para eliminar la tendencia.
  \item $D = 0$ o $1$: a determinar tras analizar la estacionalidad residual.
\end{itemize}

\paragraph{d) Modelo candidato:}
Dado el decaimiento de la ACF y el corte en la PACF,
SARIMA$(1,1,0)\times(0,0,0)_{12}$ o SARIMA$(1,1,1)$.

\paragraph{e) Validación:}
\begin{itemize}
  \item Análisis de residuos: test de Ljung-Box.
  \item Backtesting: RMSE o MAE en datos fuera de muestra.
\end{itemize}

% ══════════════════════════════════════════════════════════════════════════════
\subsection{Pregunta 6 -- Serie de Ventas Mensuales}

\begin{lstlisting}
import pandas as pd
import matplotlib.pyplot as plt

df = pd.read_csv('ventas_empresa.csv',
                 parse_dates=['Fecha'], index_col='Fecha')
ts = df['Ventas'].resample('MS').sum()

plt.figure(figsize=(10, 4))
ts.plot(title='Serie Mensual de Ventas')
plt.show()
\end{lstlisting}

\paragraph{b) Características observadas:}
\begin{itemize}
  \item \textbf{Tendencia}: creciente y persistente.
  \item \textbf{Estacionalidad}: fuerte y estable, con ciclos anuales.
  \item \textbf{No estacionariedad}: en media (tendencia) y eventualmente
        en varianza.
\end{itemize}

\paragraph{c) Diferenciación propuesta:}
$s=12$, $d=1$, $D=1$.

\paragraph{d) Modelo candidato:}
SARIMA$(1,1,0)\times(0,1,1)_{12}$.

\paragraph{e) Métricas de evaluación:}
\begin{itemize}
  \item \textbf{MAPE}: error porcentual medio absoluto; útil para interpretación
        gerencial.
  \item \textbf{RMSE}: penaliza errores grandes; recomendable cuando se desea
        evitar pronósticos con alta desviación.
\end{itemize}

% ══════════════════════════════════════════════════════════════════════════════
\subsection{Pregunta 7 -- ARIMA vs.\ ARIMAX vs.\ SARIMA vs.\ SARIMAX}

\paragraph{a) Diferencias:}
\begin{itemize}
  \item \textbf{ARIMA}: modelo univariado para series no estacionales.
  \item \textbf{ARIMAX}: ARIMA con variables exógenas.
  \item \textbf{SARIMA}: ARIMA con componentes estacionales.
  \item \textbf{SARIMAX}: versión más completa; incorpora estacionalidad y
        variables exógenas.
\end{itemize}

\paragraph{b) Ecuación general del SARIMAX:}
\[
  \Phi_P(B^s)\,\phi_p(B)\,(1-B)^d\,(1-B^s)^D X_t
  = \beta^\top Z_t + \Theta_Q(B^s)\,\theta_q(B)\,\varepsilon_t
\]
donde $\beta^\top Z_t$ es el componente de regresión exógena.

\paragraph{c) Justificación de variables exógenas:}
Se incluye una variable exógena (p.\,ej.\ gasto en publicidad o tipo de cambio)
si existe causalidad teórica, mejora de precisión observada y capacidad de
pronóstico de dicha variable.

\paragraph{d) Modelo SARIMAX candidato:}
SARIMAX$(1,1,1)\times(0,1,1)_{12}$.

\begin{lstlisting}
from statsmodels.tsa.statespace.sarimax import SARIMAX

# 'ventas' es la serie y 'exog_data' es un DataFrame con variables externas
model = SARIMAX(ventas,
                order=(1, 1, 1),
                seasonal_order=(0, 1, 1, 12),
                exog=exog_data)

results = model.fit()
print(results.summary())
\end{lstlisting}

\paragraph{f) Ventajas y limitaciones del SARIMAX frente al SARIMA:}
\begin{itemize}
  \item \textbf{Ventajas}: mayor capacidad explicativa; reduce el error de
        pronóstico cuando la variable exógena es relevante.
  \item \textbf{Limitaciones}: requiere pronosticar la variable exógena; mayor
        riesgo de sobreajuste si se incluyen variables sin relación causal genuina.
\end{itemize}

\end{document}
